\documentclass[12pt]{article}

% Title information
\title{Literature for Info Treatments}
\author{Tim Kircali, Stefanie Stantcheva and Matthias Weber}
\date{\today} % Use \date{} for no date

% Bibliography package
\usepackage[backend=biber, style=apa]{biblatex}
\addbibresource{references.bib} % Link to the bibliography file

\begin{document}

\maketitle

% Table of contents
\tableofcontents
\section{Treatment Design}
Prepost design: elicit dependent variables before and after the treatment.
"eliciting posterior first-stage variables (after the treatment) is advisable. It is the only way to get an actual first stage of your treatment and see if it worked and, if yes, by how much." --> It helps to evaluate heterogenous treatment effects. \\
"Dependent second-stage variables always have to be measured (at least) posttreatment. The question is whether you need to measure them pretreatment too. Pretreatment measures can add precision since they allow you to filter out an individual fixed effect and the tendency of a respondent to respond in a given way. As a first pass, a substitute is, of course, to control for covariates (measured pretreatment only ideally) that are relevant."

"Yet, for both first- and second-stage variables, there are concerns about asking twice: (a) Consistency pressures may prompt the respondents to answer similarly pre- and posttreatment. Yet, it is possible to ask similar, but not identical, questions to avoid making the inconsistency salient and to vary the question format. (b) Asking twice about the same thing may also lead to more EDE or SDB if it makes the respondent more aware of the topic of the experiment. (c) Related to this, the biggest worry is perhaps that, by asking these questions, you are already priming respondents to think about the topic, and priming in itself may not be neutral. (d) More elicitations also mean more time commitment for the respondents, increasing survey fatigue and cognitive load and potentially leading to the problems discussed in Supplemental Appendix A-2. The trade-off will generally depend on the topic and on how sensitive and prone to priming, EDE, and SDB it is."

\section{Info Treatments}
\subsection*{Data Privacy Treatment}
In the first information treatment, we illustrate tax-relevant players and institutions and how tax-relevant data is being exchanged between them. Then we explain how information exchange would change if AEI is expanded. For that, we consider using comparisons to other countries and explain how it works there. The primary goal of the data privacy treatment is to shock people’s knowledge and perceptions about AEI, specifically addressing their data privacy considerations and distinguishing them from other economic considerations. 
[We could also discuss withholding tax here] \\[0.5em]
\textbf{Information Text}
\begin{itemize}
    \item People usually file tax returns to communicate to the government how much income they've had.
    \item Based on these reports, the tax administration/IRS then evaluates what taxes are being charged to the taxpayer.  
    \item But the tax authorities do not only rely on the self-reported information. To check that the self reported information is correct they also use information that is third-party reported 
    \item Third-party reporting means that certain institutions are obliged to notify the tax administration when they issue payments that might be tax relevant. That means they automatically exchange data with the tax autorities.
    \item For example, employers report to the IRS how much they pay their employees in wages. 
    \item Another example is that if interest and dividend payments of assets are higher than a set amount, the tax authority will also be notified automatically. 
    \item (Additional information on which data is being automatically exchanged)
    \item If information about certain tax payments is not automatically exchanged, it is hard for the tax authorities to check that the reported income is correct. 
    \item In these cases the tax administration can either rely on the information provided by the taxpayer or they can take investigative and legal steps to uncover more information, while the latter can be time consuming and costly for the tax administration.
    \item The level of third-party reporting, that is the amount of different data sources that is automatically reported to the tax authority, differs among countries. While in some countries nearly everything, from employer data to bank account and health insurance data, is reported, in other countries tax authorities have no automatic access to tax relevant data at all.   
    \item Although automatic information reporting helps tax administrations directly verify the accuracy of tax reports, some people are concerned about their financial data privacy. They don't like the fact that governmental institutions collect and exchange private data about them and think it is important that the amount of data that is being exchanged with the government should be kept at a minimum.
    \item Still, one must consider that eventually the government will end up having most of the data of honest taxpayers because they will give it to them. They just lack of access of data of dishonest taxpayers.  
\end{itemize}
\textbf{Other Details}
\begin{itemize}
    \item According to chatgpt this information treatment would last 2.5 to 3 minutes.
    \item Detailed information about which data types that are exactly exchanged are not really important and the results would also not be applicable to many contexts. That's why we keep the information about automatic exchange of information rather general.  
\end{itemize}

\subsection*{Tax Filing Treatment}
The second treatment discusses concerns about tax filing costs such as the fact that an average American taxpayer has 12 hours + 270 \$ tax filing costs. It describes how (guided) tax software works and what a prepopulated tax return looks like. We emphasize that tax software and prepopulated tax returns help tax filers to make fewer mistakes and significantly reduce their tax filing costs. And, most importantly, we highlight the importance of AEI to the tax administration in providing prepopulated tax returns. The primary goal of the tax filing treatment is to shock people's knowledge and perceptions about tax filing costs and tax filing services and how they are related to the existence of AEI.  \\[0.5em]
\textbf{Information Text}
\begin{itemize}
    \item People usually file tax returns to tell the government how much income they've earned.
    \item While some taxpayers do their taxes themselves, others pay for tax advice or pay/request someone else to prepare their tax returns. 
    \item Another option for taxpayers is to use tax software to file their taxes. These help taxpayers file their tax returns by providing guidance through the tax formulas and explaining which deductions might apply and which ones do not.  It gives field-by-field tips and tricks to reduce errors and foregone reductions.
    \item All these options are either time consuming or costly. That is why taxpayers face considerable tax filing costs each year. An estimate by the IRS suggests that an average taxpayer in the US has 12 hours + 270\$ tax filing costs each year.
    \item Some countries' tax administration offer solutions that reduce tax filing costs. However, the extend to which they do so differs among countries. While most taxpayers in the US have to pay for tax software licences, many other countries offer free government provided tax software solutions to reduce tax filing costs. 
    \item Some countries even offer prepopulated tax returns, which is a more advances version of tax software. This means that the tax returns are pre-filled with information the tax authority already has about the taxpayer. This makes filing your tax return much quicker and less likely to have mistakes. 
    \item In order to offer prepopulated tax returns, it is very important that the government has detailed financial information about the taxpayer. Otherwise, the tax administration will not have enough information to prefill the returns.
    \item An important factor that helps tax administrations to get access to detailed financial information is third-party reporting and the automatic exchange of information of institutions with tax authority in general. That means that certain institutions are obliged to notify the tax administration about payments and other information that might be tax relevant. 
    \item The more information the tax administration automatically receives, the better it can provide pre-populating services. It's not surprising that only countries that automatically exchange information with lots of institutions offer services that fill in forms for people.
    \item Even if government provided tax filing services such as free tax software and prepopulated tax returns bring many advantages, there are also people that argue against governments efforts to provide them.
    \item For example, they argue that tax software is already being offered by commercial providers and that government provided software might not reach its quality. 
    \item Another argument is that the IRS should not act as both the tax collector and tax preparer. This could lead to conflicts of interest or undermine the interest representation of taxpayers.
\end{itemize}
\\
\textbf{Other Details}
\begin{itemize}
    \item According a video with that text would last between 3 to 4 minutes. 
\end{itemize}

\subsection*{Tax Evasion Treatment}
The third treatment provides respondents with information about the economic effects of AEI in tax enforcement. We give them information about the size of the tax gap, about the limited resources of IRS and that audit levels are generally low. We directly point out that AEI is effective in reducing tax evasion by explaining evidence which shows that tax evasion is rare in areas which are covered under AEI. The primary goal of the tax evasion treatment is to shock people’s knowledge and perceptions about the economic effects of tax enforcement and AEI without affecting considerations related to fairness \& inequality. \\[0.5em]
\textbf{Information Text}
\begin{itemize}
\item Tax revenue is essential for modern governments to provide their functions, reaching from the provision of safety over infrastructure, and education to healthcare and social security.
\item One problem governments face is the loss of tax revenue due to tax evasion. This means that some taxpayers deliberately hide some of their income sources from the tax authorities and the government. 
\item The tax gap is the difference between the taxes that should be paid and the taxes that are actually paid. Recent estimates show that this gap is shockingly high in many countries.
\item In the US, for example, the IRS's latest estimates suggest that only 85\% of the total actual tax is paid voluntarily and on time. Of the remaining 15\%, which amounts to \$696 billion, about \$90 billion is either paid voluntarily but late or collected through IRS administrative and enforcement activities. Therefore, 13.1\% of the total actual tax, or \$606 billion, is ultimately not collected by the IRS.
\item (Information about the situation in other countries)
\item If the government was able to collect these taxes, it could use them for a number of things. To give a few examples, it could double spending on education (\$677 billion), it could increase social security (\$1,219 billion) by 50\%, or the government could simply reduce taxes for all taxpayers by this 13.1\%. 
\item Tax administrations have a number of options to combat tax evasion and reduce the tax gap.
\item Recent research has shown that an increase in audit activity leads to an increase in tax revenue that far outweighs the additional costs of auditing.
\item Another important factor that helps tax administrations to reduce tax evasion is third party reporting and automatic exchange of information with the tax authority in general. This means that certain institutions are obliged to inform the tax administration about payments and other information that may be relevant for tax purposes.
\item Research has shown that tax evasion takes place almost exclusively in areas not covered by automatic information exchange. Increasing the amount of data covered by the automatic exchange of information is therefore an effective way of reducing tax evasion.
\end{itemize}
\\
\textbf{Other Details}
\begin{itemize}
    \item The detailed numbers for the tax gap come from the IRS report "Tax Gap Projections for Tax Year 2022".
    \item These numbers refer to the federal taxes that account for around 2/3 of overall taxes in the US. 
    \item The information about the public spending come from wikipedia \url{https://en.wikipedia.org/wiki/2022_United_States_federal_budget#/media/File:Federal_budget_2022.webp}
    \item We use information of Boning et al about the effectiveness of tax audits
    \item We use information about the effectiveness of third party reporting from literature (Kleven et al., Pomeranz et al. and so on)
    \item reading time 3 to 3.5 minutes according to chatgpt
\end{itemize}

\subsection*{Inequality Treatment}
The last specific information treatment again includes information on tax evasion, but also explains the fact that the tax gap is much more pronounced for high-income individuals, and discusses the reasons why this is the case. After that, we explain fairness considerations of tax evasion such as the fact that honest taxpayers have to finance the fiscal gap caused by tax evaders. The primary goal of the inequality treatment is to shock people’s knowledge and perceptions about fairness and distributional considerations of tax enforcement policy and AEI. 
\textbf{Information Text}
\begin{itemize}
    \item The tax gap is the difference between the taxes that should be paid and the taxes that are actually paid. Recent estimates show that this gap is shockingly high in many countries.
    \item In the US, for example, the IRS's latest estimates suggest that only 85\% of the total actual tax is paid voluntarily and on time. Of the remaining 15\%, which amounts to \$696 billion, about \$90 billion is either paid voluntarily but late or collected through IRS administrative and enforcement activities. Therefore, 13.1\% of the total actual tax, or \$606 billion, is ultimately not collected by the IRS.
    \item (Here we can give information about how tax gaps are distributed along the income distribution, for example we could use examples Johns and Slemrod (2010))
    \item We can give information that most of the misreporting happens in understated income accounts and that most of them are result of misreported business income
    \item (We can discuss reasons why this is the case )
    \item (And we could discuss the fact that honest taxpayers finance the tax gap created by dishonest taxpayers)    
    \item Recent research has shown that an increase in audit activity leads to an increase in tax revenue that far outweighs the additional costs of auditing.
    \item Here we can additionally inform them that returns to audits are much higher when auditing higher income individuals
    \item 
    \item Another important factor that helps tax administrations to reduce tax evasion is third party reporting and automatic exchange of information with the tax authority in general. This means that certain institutions are obliged to inform the tax administration about payments and other information that may be relevant for tax purposes.
    \item Research has shown that tax evasion takes place almost exclusively in areas not covered by automatic information exchange. Increasing the amount of data covered by the automatic exchange of information is therefore an effective way of reducing tax evasion.

\end{itemize}
\textbf{Other Details}
\begin{itemize}
    \item As discussed in Slemrod (2007), about two-thirds of all underreporting of income happens on the individual income tax. Understated income accounts for over 80 percent of individual underreporting of tax. Business income, as opposed to wages or investment income, accounts for about two thirds of the understated individual income.
    \item We can use estimates of Johns and Slemrod (2010) about the misreporting percentages (I think they have it form the IRS). (Table 2) They also show estimates how much taxes each individual is saving in percentage of actual taxes they should pay. (These results show a different angle because even if misreporting is much lower for low incomes, that translates much stronger into saved taxes because marginal taxes are very high at low incomes). 

\end{itemize}




\subsection*{Full Info Treatment}
The full information treatment will include all the information above and will therefore be much longer than the specific info treatments. We plan to implement a test that addresses the issues regarding survey fatigue. And finally, the control group will not receive information. 

\section{Implementation}
Where can we create these videos? Vyond, Powtoon, Canva, PowerPoint, Clip 


\subsection{AEI Info treatments}
\textbf{Info about tax gaps (together):} Here we should have some info from the IRS about the situation in the US. Possibly we could also include other countries. The question is a bit if we also give them information about tax gaps across the income distribution or other subgroups. \\
\textbf{Info about effects on evasion and tax revenues after introducing AEI:} There is not much literature about the effects on tax evasion if automatic exchange is introduced. Rather there is literature about the effects that no evasion is detected when we have an audit shock. ((\cite{kleven2011unwilling}),(\cite{pomeranz2015no}))    
\textbf{Info about which data the government \& authorities have access to:} Here one option would be to provide different information to people of different countries. Another option would be to give the same information to all population groups. Somehow also providing them with information about how other countries are enforcing. (If we take this option we could easily bring the bridge to the prefilled tax declarations)\\
\textbf{Info about pre-filled tax declarations and reduction in tax filing costs and efforts:} We can explain (on a country example such as Germany) that when a country has automatic information on many tax relevant stuff they reduce effort to control tax declarations but they can also provide pre-filled tax declarations. 
\subsection{Tax Enforcement Info Treatment}
\textbf{Info about tax gaps (together):} Here we should have some info from the IRS about the situation in the US. Possibly we could also include other countries. The question is a bit if we also give them information about tax gaps across the income distribution. \\
\textbf{Info about effects on evasion and tax revenues after increasing tax enforcement investments:} We could use the estimates of \cite{kleven2011unwilling} to explain return to audits. \\
\textbf{Info about punishment and fines:} This is probably very hard to implement as rules might be very different across countries!
\subsection{Inequality Treatment}


\section{relevant literature on tax gaps and evasion distribution}
\cite{guyton2021tax}: Tax evasion at the top of the income distribution: Theory and Evidence. 
Find that the IRS's Random audit strategy strongly underestimates evasion rates in the top income bracket , because the observed evasion on small subsamples are already surpassing predictions by IRS's random audit strategies. They then evaluate the role of offshore evasion and evasion on pass-through business income



\section{relevant literature on returns to audits}
\cite{boning2025welfare}: "We estimate the returns to IRS audits of taxpayers across the income distribution. We find an additional \$1 spent auditing taxpayers above the 90th income percentile yields more than \$12 in revenue, while audits of below-median income taxpayers yield \$5." These measures include initial audit effects but also deterrence effects which are relatively more important --> More details in paper. in contrast to other papers they also calculate indirect costs of increasing tax audits. From that paper we could also use estimates on how much it cost to audit taxpayers, but also on. 

\cite{debacker2018once}: "We examine the impact of enforcement on subsequent compliance behavior by taxpayers. Exploiting waves of randomized audits ... we find three long-run responses. First, audits increase tax payments substantially in following years, but this effect is short-lived when third-party reporting is not available. Second, taxpayers with high income volatility revert to their preaudit behavior quickly. Third, sophisticated taxpayers are affected less by enforcement."  
We find that an audit increases subsequent reported wage income by 1.3 percent and sole-proprietorship income, reported on schedule C of form 1040, by 14.2 percent on average. Thus, the effect of audit is much larger when there is less third-party reporting, which is consistent with Kleven et al. (2011).

\section{Literature on returns to AEI}: 
\cite{debacker2018once}: We find that an audit increases subsequent reported wage income by 1.3 percent and sole-proprietorship income, reported on schedule C of form 1040, by 14.2 percent on average. Thus, the effect of audit is much larger when there is less third-party reporting, which is consistent with Kleven et al. (2011).

\cite{kleven2011unwilling}: We find a very strong variation in tax evasion depending on the information environment. For third-party reported income, the evasion rate is always extremely small: it is equal to 0.23\% for total positive income, 0.35\% for total negative income, and always below 1\% across all the different categories. Interestingly, the evasion rate for self-employment income conditional on third-party reporting is only 0.33\%, suggesting that overall tax evasion among the self-employed is large because of the information environment and not because of, for example, different preferences among those choosing self-employment (such as attitudes toward risk and cheating). By contrast, tax evasion for self-reported income is substantial: the evasion rate is 17.1\% for total positive income, 7.5\% for total negative income, 5.4\% for capital income, 13.6\% for stock income, and 17.7\% for self-employment income.
"First, the probability of evading jumps up immediately once the taxpayer has some income that is self-reported (although it never exceeds 40\%). Second, the share of total income evaded is increasing in the share of income that is self-reported, whereas the share of third-party income evaded is always very close to zero. This shows that taxpayers with more self-reported income evade more, but always declare third-party income fully."


This is an example citation \cite{boning2025welfare}.

\cite{baselgia2023compliance}: Finds that tax evasion might be much more distributed in Switzerland, also among the middle class. And that perceived thread to AEoI has affected tax compliance. 




% Print bibliography
\printbibliography[heading=bibintoc]

\end{document}